\documentclass[11pt,a4paper]{article}
\usepackage[utf8]{inputenc}
\usepackage[T1]{fontenc}
\usepackage{lmodern}
\usepackage{geometry}
\usepackage{hyperref}
\usepackage{longtable}
\usepackage{array}
\usepackage{enumitem}
\geometry{margin=1in}

\title{Socratic AI Tutor: Full Project Documentation}
\author{Artificial Intelligent Virtual Assistant Project}
\date{\today}

\begin{document}
\maketitle
\tableofcontents
\newpage

\section{Project Overview}
This project is an educational virtual assistant that combines document retrieval with language model generation (RAG: Retrieval-Augmented Generation). The system allows users to upload study files, store them as vector embeddings, and ask grounded questions about those files.

\subsection{Main Goals}
\begin{itemize}[leftmargin=*]
\item Provide accurate, evidence-grounded answers from user-provided files.
\item Keep memory persistent across restarts.
\item Allow users to manage memory: save files, delete one file, or reset all.
\item Offer built-in quality assessment for retrieval and answer generation.
\end{itemize}

\section{System Architecture}
The app is structured as three UI panels:
\begin{itemize}[leftmargin=*]
\item Left: Knowledge base management (upload/save/list/delete/reset).
\item Center: Chat interface (user interaction with tutor).
\item Right: Assessment and evaluation results.
\end{itemize}

\subsection{Core Components}
\begin{longtable}{|p{0.28\linewidth}|p{0.65\linewidth}|}
\hline
\textbf{File} & \textbf{Responsibility} \\
\hline
\texttt{app.py} & Streamlit UI, user workflow, chat loop, panel layout, button logic. \\
\hline
\texttt{database.py} & File parsing, chunking, embedding ingestion, retrieval helpers, file-level delete/reset operations. \\
\hline
\texttt{tutor_engine.py} & Response logic, chitchat detection, grounding checks, context handling, fallback behavior. \\
\hline
\texttt{evaluation.py} & Golden dataset loading and metric calculation (Hit@k, MRR, token F1, semantic similarity). \\
\hline
\texttt{config.py} & Runtime constants and environment variable mapping. \\
\hline
\texttt{docker-compose.yml}, \texttt{Dockerfile} & Containerized runtime and reproducible deployment. \\
\hline
\end{longtable}

\section{Data and Memory Flow}
\subsection{Ingestion Pipeline}
\begin{enumerate}[leftmargin=*]
\item User uploads a file (\texttt{txt, md, pdf, csv, json}).
\item The file is converted to clean text.
\item Text is chunked using recursive character splitter.
\item Chunks are embedded and stored in Chroma persistent vector storage.
\item Stable IDs prevent duplicate chunk insertion for the same source/chunk index.
\end{enumerate}

\subsection{Retrieval and Answer Pipeline}
\begin{enumerate}[leftmargin=*]
\item User sends a chat message.
\item If message is chitchat (greeting/thanks/goodbye), answer directly (no retrieval).
\item Otherwise retrieve top-k chunks (MMR) and build contextual prompt.
\item Apply grounding checks and concise answer strategy.
\item Return answer and evidence metadata (source/page/chunk).
\end{enumerate}

\section{Supported File Types}
\begin{itemize}[leftmargin=*]
\item \texttt{.txt}, \texttt{.md}: direct text read.
\item \texttt{.pdf}: extracted with PyPDF2 page text extraction.
\item \texttt{.csv}: rows flattened to text lines.
\item \texttt{.json}: nested objects flattened to key-value lines.
\end{itemize}

\section{Persistence Model}
\subsection{Expected Behavior}
\begin{itemize}[leftmargin=*]
\item First run starts empty.
\item Saved knowledge persists between runs.
\item Users can delete a single file or clear all knowledge.
\end{itemize}

\subsection{Docker Volumes}
The project uses named Docker volumes for app memory folders. This avoids accidental preloaded content from local project folders and keeps runtime memory explicitly managed by app actions.

\section{Evaluation Framework}
\subsection{Golden Dataset}
Evaluation uses a JSONL dataset with:
\begin{itemize}[leftmargin=*]
\item question
\item reference answer
\item expected source file
\end{itemize}

\subsection{Metrics}
\begin{itemize}[leftmargin=*]
\item Hit@k: whether expected source appears within top-k retrieval.
\item MRR: ranking quality of expected source.
\item Token F1: lexical overlap between generated and reference answers.
\item Semantic Similarity: embedding cosine similarity between generated and reference answers.
\end{itemize}

\section{UI/UX Decisions}
\begin{itemize}[leftmargin=*]
\item Chat input is centered with chat panel to improve visual consistency.
\item Per-file delete button is compact for cleaner saved-file list display.
\item \texttt{Evaluate Model} remains disabled until at least one file is saved.
\item Result panel is shown in expander to avoid popover layering issues.
\end{itemize}

\section{Running the Project}
\subsection{Docker (Recommended)}
\begin{enumerate}[leftmargin=*]
\item Create \texttt{.env} with valid OpenAI key and model.
\item Run:
\begin{verbatim}
docker compose up -d --build
\end{verbatim}
\item Open \texttt{http://localhost:8501}
\item Stop with:
\begin{verbatim}
docker compose down
\end{verbatim}
\end{enumerate}

\subsection{Local Python}
\begin{enumerate}[leftmargin=*]
\item Create and activate virtual environment.
\item Install dependencies:
\begin{verbatim}
pip install -r requirements.txt
\end{verbatim}
\item Run:
\begin{verbatim}
streamlit run app.py
\end{verbatim}
\end{enumerate}

\section{Configuration}
Main configuration values are in \texttt{config.py}:
\begin{itemize}[leftmargin=*]
\item chunk size and overlap
\item default retrieval k
\item embedding model name
\item OpenAI model/temperature/max output tokens
\end{itemize}

\section{Operational Checklist}
\begin{enumerate}[leftmargin=*]
\item Upload file and save.
\item Verify file appears under Saved Files.
\item Ask question and verify grounded answer.
\item Run evaluation and inspect metrics.
\item Delete one file and confirm it disappears.
\item Reset knowledge base and confirm evaluation becomes disabled.
\end{enumerate}

\section{Known Limitations and Next Steps}
\begin{itemize}[leftmargin=*]
\item PDF extraction quality depends on PDF text layer (scanned PDFs may require OCR).
\item Current semantic evaluation is embedding-based; optional LLM-judge can be added.
\item Deprecation warnings in LangChain APIs can be addressed by future package upgrades.
\item Additional file formats (e.g., DOCX) can be added in ingestion module.
\end{itemize}

\section{Conclusion}
The current system provides a robust educational RAG assistant with persistent memory, transparent retrieval evidence, file-level memory controls, and reproducible deployment. The architecture is modular and ready for iterative improvement in model quality, evaluation depth, and file-format support.

\end{document}
